\selectlanguage{ngerman}%
\frenchspacing%
\begin{ledgroupsized}[r]{120mm}%
\footnotesize%
\pstart%
\noindent%
\textbf{Überlieferung:}%
\pend%
\end{ledgroupsized}%
\begin{ledgroupsized}[r]{114mm}%
\footnotesize%
\pstart%
\parindent -6mm%
\makebox[6mm][l]{\textit{L}}%
Aufzeichnung:
LH XXXVIII Bl.\;91.
Ein Blatt, 8\textsuperscript{o};
linker, oberer und rechter Rand beschnitten;
Fragment eines \href{\Scan{LH_38_091wz}}{Wassezeichens}:
Niedersächsiches Papier.
Eine Seite auf Bl.\;\href{\Scan{LH_38_091r}}{91\,r\textsuperscript{o}}.
Bl.\;\href{\Scan{LH_38_091v}}{91\,v\textsuperscript{o}} ist unbeschrieben.
\pend%
\end{ledgroupsized}%
\vspace*{5mm}%
\begin{ledgroup}%
\footnotesize%
\pstart%
\noindent%
\textbf{Datierungsgründe:}
Die Aufzeichnung N.\;\ref{55774} ist von Leibniz eigenhändig datiert.
Ihm gegenüber erwähnte Jean Paul de la Roque%
\indexN{\textls{La Roque} (Larroque), Jean Paul de, gest. 1691}
in seinem Brief vom 3.\;Juni 1678 einen zeitgenössischen Versuch,
die \textit{Colombe artificielle d'Architas qui voloit ... une trentaine de pas}
nachzuahmen (\textit{LSB} III,\;2 N.\;175, S.\;456.15--17).\cite{01477}
Dies dürfte Leibniz angespornt haben, die Aufzeichnung N.\;\ref{55774} zu verfassen.
\pend%
\end{ledgroup}%
\selectlanguage{latin}%
\frenchspacing%
\vspace{8mm}%
\normalsize%
\pstart%
\noindent%
\lbrack%
\href{\Scan{LH_38_091r}}%
{91\,r\textsuperscript{o}}%
\rbrack%
\phantom{Junii 1678}%
\hfill%
\phantom{\lbrack91\,r\textsuperscript{o}\rbrack}%
Junii 1678
\pend%
\pstart%
\noindent%
\centering%
\edtext{%
Columba Architae%
}{%
\lemma{Columba Architae}%
\Cfootnote{%
Von Archytas' mechanischer Taube, die angeblich zu fliegen vermochte, berichtet
\textsc{Aulus Gellius}, \textit{Noctes Atticae} X,\;12,\;9--10.\cite{01478}
}}
resuscitata%
\indexN{\textls{Archytas} von Tarent (um 365 v. Chr.)}%
\indexS{columba Architae}
\pend%
\vspace{0.5em}%
\pstart%
\noindent%
Quaeritur magna vis%
\indexS{vis magna inclusa moli exiguae}%
\indexS{moles exigua}
\edtext{%
exigua in mole.%
}{%
\lemma{exigua in}%
\Bfootnote{%
\textit{(1)}~spatio %
\textit{(2)}~mole%
~\textit{L}%
}}
Haec vis quidem est pulveris pyrii:%
\indexS{pulvis pyrius}%
\indexS{vis pulveris pyrii}
sed in ea~duo incommoda,%
\indexS{incommodum}
primo,
quod tardari non
\edtext{%
potest,
\lbrack nec\rbrack\
machinis%
}{%
\lemma{potest,}%
\Bfootnote{%
\textit{(1)}~sed %
\textit{(2)}~\bars{et \textit{ändert Hrsg.}} machinis%
~\textit{L}%
}}
accommodari,
sed momento absolvit ictum;%
\indexS{ictus pulveris pyrii}
ideoque pro alarum columbae ictibus%
\indexS{ictus alae}%
\indexS{ala}%
\indexS{columba Architae}
totidem opus esset
pulveris displosionibus;%
\indexS{displosio pulveris pyrii}%
\indexS{pulvis pyrius}
alterum incommodum est%
\indexS{incommodum}
quod se prodit strepitu,%
\indexS{strepitus pulveris pyrii}
fumo,%
\indexS{fumus pulveris pyrii}
flamma,%
\indexS{flamma pulveris pyrii}
\edtext{%
unde parum gratus%
}{%
\lemma{unde}%
\Bfootnote{%
\textit{(1)}~grati\textlangle am\textrangle\ %
\textit{(2)}~parum gratus%
~\textit{L}%
}}
erit volatus,%
\indexS{volatus}
et similis videbitur racquetis%
\indexS{raquetum}
sive
\edtext{%
fusis ignium%
}{%
\lemma{fusis}%
\Bfootnote{%
\textit{(1)}~ignorum %
\textit{(2)}~ignium%
~\textit{L}%
}}
artificialium.%
\indexS{fusus ignum artificialium}%
\indexS{ignis artificialis}
Vera
\edtext{%
ratio exiguae%
}{%
\lemma{ratio}%
\Bfootnote{%
\textit{(1)}~exiguam v %
\textit{(2)}~exiguam %
\textit{(3)}~exiguae%
~\textit{L}%
}}
moli%
\indexS{moles exigua}
magnam vim includendi,%
\indexS{vis magna inclusa moli exiguae}%
\indexS{ratio includendi vim magnam moli exiguae}
est per pressionem aeris.%
\indexS{pressio aeris}%
\indexS{aer compressus}
\edtext{%
Nam aer in decimam%
}{%
\lemma{Nam}%
\Bfootnote{%
\textit{(1)}~aer s %
\textit{(2)}~aer in %
\textit{(a)}~se\textlangle--\textrangle\ %
\textit{(b)}~decimam%
~\textit{L}%
}}
spatii partem redactus%
\indexS{aer redactus in partem spatii}
tantam vim habebit,%
\indexS{vis aeris compressi}
quantam columna aquae%
\indexS{columna aquae}
300 pedum%
\indexS{pes}
ejusdem basis%
\lbrack;\rbrack\
quod pondus in
\edtext{%
quantulacunque latitudine%
}{%
\lemma{quantulacunque}%
\Bfootnote{%
\textit{(1)}~altitudine %
\textit{(2)}~latitudine%
~\textit{L}%
}}
immensum est.%
\indexS{pondus immensum}
Hinc facile erit elaborare
\edtext{%
horologium portatile%
\indexS{horologium portatile}
unius anni%
\indexS{annus}%
}{%
\lemma{horologium}%
\Bfootnote{%
\textit{(1)}~unius ann %
\textit{(2)}~intra %
\textit{(3)}~portatile unius anni%
~\textit{L}%
}}
et ultra,
imo perpetuum%
\indexS{horologium perpetuum}%
\lbrack,\rbrack\
\edtext{%
de quo alibi.%
}{%
\lemma{de quo alibi}%
\Cfootnote{%
Wohl Anspielung auf N.\;\ref{motus-perpetuus_intro} (Oktober 1677) in diesem Band.
}}
Nunc mediocri vi aeris,%
\indexS{vis aeris compressi}
inclusa
\edtext{%
intra speciem%
}{%
\lemma{intra}%
\Bfootnote{%
\textit{(1)}~modum %
\textit{(2)}~speciem%
~\textit{L}%
}}
follis%
\indexS{fluidum aperiens se}
vel parallelepipedi se%
\indexS{parallelepipedum aperiens se}
\edtext{%
aperientis
et
\lbrack se\rbrack\
aperiendo
alas%
\indexS{ala}
latas sed tenues columbae%
\indexS{columba Architae}%
}{%
\lemma{aperientis}%
\Bfootnote{%
\bars{et claudentis, \textit{gestr.}} %
et %
\bars{se \textit{erg.\;Hrsg.}} aperiendo alas %
\bars{%
\textit{(1)}~latas sed tenues %
\textit{(2)}~latas sed tenues %
\textit{erg.}} columbae%
~\textit{L}%
}}
ex levissima materia%
\indexS{materia levissima}
confectae
\edtext{%
fortiter agitantis,%
}{%
\lemma{fortiter}%
\Bfootnote{%
\textit{(1)}~circumagentis, %
\textit{(2)}~agitantis,%
~\textit{L}%
}}
aer externus%
\indexS{aer externus}
percutietur%
\indexS{aer percussus ab ala}
et columba elevabitur,%
\indexS{columba Architae}
repetitis per intervalla alarum ictibus.%
\indexS{ictus alae}%
\indexS{ictus repetitus per intevalla}%
\indexS{intervallum}
Effici etiam potest,
ut se huc illuc flectat,
quasi ope temonis.%
\indexS{temo}
Quin et effici potest,
ut levissima mutatio caloris%
\indexS{mutatio caloris}%
\indexS{calor mutatus}
et frigoris%
\indexS{mutatio frigoris}%
\indexS{frigus mutatum}
denuo ipsum animet aerem
comprimendo%
\indexS{compressio aeris}%
\indexS{aer compressus}
vel dilatando,%
\indexS{dilatatio aeris}%
\indexS{aer dilatatus}
item
ut columba%
\indexS{columba Architae}
decidens
ipso lapsus ictu%
\indexS{ictus lapsus}%
\indexS{lapsus}
denuo aerem comprimat,%
\indexS{aer compressus}
eo artificio perfacili%
\indexS{artificium perfacile}
ut comprimi%
\indexS{compressio aeris}
\edtext{%
possit momento,
non momento dilatari.%
\indexS{dilatatio aeris}%
}{%
\lemma{possit}%
\Bfootnote{%
\textit{(1)}~, dilatare se %
\textit{(2)}~momento, non momento dilatari.%
~\textit{L}%
}}
Sic post casum%
\indexS{casus}
se denuo elevabit ad volatum.%
\indexS{volatus}
\pend%
\selectlanguage{latin}%
\frenchspacing%
\normalsize%
